\documentclass[a4paper]{article}
\usepackage{graphicx}
\usepackage{amsmath,amssymb}
\usepackage{hyperref}
\usepackage{minted}
\usepackage{multirow}

\title{Abstract}
\date{}
\begin{document}
    \maketitle
    This research proposes a framework for self-awareness in multi-robot systems by suggesting the evolution of sensory action-potential trains into a mutual language inspired by neurology and classical linguistics. Here, the highest form of self-awareness refers to awareness of the robot's own awareness, in the sense that the system can use this awareness for autonomous navigation and for updating its generative model. By language, this research refers to Saussure’s syntagmatic and paradigmatic relations and Lévi-Strauss’s theory of binary oppositions. The framework first builds an internal language based on event-specific knowledge gained through guided learning. It then shares its essential constituents with the other robots to distinguish self from other and novel sensory data from what has already been experienced and memorized.
\end{document}
